\documentclass[11pt]{article}
\usepackage[margin=1in]{geometry}
\usepackage{amsmath,amssymb,amsthm}
\usepackage{booktabs}
\usepackage{float}
\usepackage{enumitem}
\usepackage{hyperref}
\usepackage{natbib}
\bibliographystyle{apalike}

\newtheorem{theorem}{Theorem}
\newtheorem{lemma}[theorem]{Lemma}
\newtheorem{proposition}[theorem]{Proposition}
\newtheorem{definition}[theorem]{Definition}

\title{Reward Recovery in Sequential Decision Models}
\author{Pranjal Rawat\thanks{Georgetown University. \texttt{pp712@georgetown.edu}.}}
\date{April 2026}

\begin{document}
\maketitle

%% ======================================================================
\section{Setup and Notation}
%% ======================================================================

Consider a stationary infinite-horizon MDP with finite state space $\mathcal{S}$, finite action space $\mathcal{A}$, discount factor $\beta \in (0,1)$, transition kernel $P : \mathcal{S} \times \mathcal{A} \to \Delta(\mathcal{S})$, and unknown reward $r : \mathcal{S} \times \mathcal{A} \to \mathbb{R}$. Additive T1EV shocks yield the soft Bellman recursion.
\begin{align}
Q^*(s,a) &= r(s,a) + \beta \sum_{s'} P(s' \mid s,a)\, V^*(s'), \label{eq:bellman} \\
V^*(s) &= \log \sum_{a'} \exp Q^*(s,a'), \label{eq:value} \\
\pi^*(a \mid s) &= \exp\bigl(Q^*(s,a) - V^*(s)\bigr), \label{eq:policy} \\
A^*(s,a) &= Q^*(s,a) - V^*(s) = \log \pi^*(a \mid s). \label{eq:advantage}
\end{align}
The advantage $A^*$ is identified from CCPs. The reward $r$ is not: any $\tilde{r}(s,a) = r(s,a) + \lambda(s)$ generates identical CCPs \citep{magnac2002}. Following \citet{kalouptsidi2021}, call the set of all such $\tilde{r}$ the identified set.

%% ======================================================================
\section{Four Approaches to Reward Recovery}
%% ======================================================================

\subsection{Reduced-Form Q}

The reduced-form approach estimates $Q_\omega(s,a)$ by logit MLE.
\begin{equation}\label{eq:rf}
\hat\omega = \arg\max_\omega \sum_{i,t} \log \frac{\exp Q_\omega(s_{it},a_{it})}{\sum_{a'}\exp Q_\omega(s_{it},a')}.
\end{equation}
At the population optimum, $Q_{\hat\omega}(s,a) = Q^*(s,a) + c(s)$ for some action-independent $c$, so the backed-out ``reward'' is the advantage $A^*(s,a)$, not the true reward $r(s,a)$.

\subsection{AIRL \citep{fu2018}}

The AIRL discriminator decomposes as follows.
\begin{equation}\label{eq:decomp}
f_\theta(s,a,s') = g_\theta(s,a) + \beta\, h_\phi(s') - h_\phi(s),
\end{equation}
Here $g_\theta$ approximates the reward and $\beta h(s') - h(s)$ is potential-based shaping \citep{ng1999}. At the GAN optimum the discriminator score equals the advantage, $f^* = A^*$ \citep[Appendix A.4]{fu2018}, which implies the following expression for the recovered reward network.
\begin{equation}\label{eq:gstar}
g^*(s,a) = r(s,a) - \delta(s) + \beta \sum_{s'} P(s' \mid s,a)\,\delta(s'),
\end{equation}
where $\delta(s) \equiv V^*(s) - h^*(s)$. Unless $\delta \equiv 0$, the recovered $g^*$ is a potential-based perturbation of $r$.

\subsubsection{LSW Anchor Identification \citep{lee2026}}

LSW impose two anchors. First, the exit payoff is normalized to zero, so $g(s,\text{exit}) = 0$ for all $s$. Second, the shaping function is pinned at the absorbing state, so $h(s_{\mathrm{abs}}) = 0$, where exit deterministically transitions to $s_{\mathrm{abs}}$.

\begin{lemma}[Non-identifiability]\label{lem:nonid}
If $(g,h)$ and $(\bar{g},\bar{h})$ produce the same $f$ for all $(s,a,s')$, then $\bar{h} = h + \Phi$ and $\bar{g}(s,a) = g(s,a) + \Phi(s) - \beta\Phi(s')$ for some $\Phi : \mathcal{S} \to \mathbb{R}$.
\end{lemma}
\begin{proof}
Rearrange $g + \beta h(s') - h(s) = \bar{g} + \beta\bar{h}(s') - \bar{h}(s)$ and define $\Phi = \bar{h} - h$. \qed
\end{proof}

\begin{theorem}[Anchor identification]\label{thm:anchor}
Under (A1)--(A2), if $(g,h)$ and $(\bar{g},\bar{h})$ realize the same $f$, then $g = \bar{g}$ and $h = \bar{h}$.
\end{theorem}
\begin{proof}
By Lemma~\ref{lem:nonid}, $\bar{g}(s,\text{exit}) = 0 + \Phi(s) - \beta\Phi(s_{\mathrm{abs}})$. Anchor (A1) forces $\Phi(s) = \beta\Phi(s_{\mathrm{abs}})$ for all $s$. Anchor (A2) forces $\Phi(s_{\mathrm{abs}}) = 0$. Hence $\Phi \equiv 0$. \qed
\end{proof}

\begin{theorem}[Recovery of true reward]\label{thm:recovery}
If the true pair $(r^*, V^*)$ satisfies $r^*(s,\mathrm{exit}) = 0$ and $V^*(s_{\mathrm{abs}}) = 0$, then any AIRL pair $(g,h)$ realizing $A^*$ under the same anchors satisfies $g = r^*$ and $h = V^*$.
\end{theorem}
\begin{proof}
Define $\Phi = h - V^*$. The exit anchor gives $\Phi(s) = \beta\Phi(s_{\mathrm{abs}})$. The absorbing anchor applied to both $h$ and $V^*$ gives $\Phi(s_{\mathrm{abs}}) = 0$. \qed
\end{proof}

The \citet{kang2025} ERM framework uses the same anchor principle (fixing the reward of a known action) but solves a single non-adversarial optimization.

\subsection{IQ-Learn \citep{garg2021}}

The inverse Bellman operator $r_Q(s,a) \equiv Q(s,a) - \beta \sum_{s'} P(s' \mid s,a)\, V_Q(s')$ expresses reward as a function of $Q$ alone. IQ-Learn estimates $Q$ by solving the following optimization problem.
\begin{equation}\label{eq:iq_obj}
\max_Q \;\mathbb{E}_{(s,a) \sim \rho^*}\bigl[\phi\bigl(r_Q(s,a)\bigr)\bigr] - (1-\beta)\,\mathbb{E}_{s_0 \sim \mu_0}\bigl[V_Q(s_0)\bigr],
\end{equation}
The function $\phi(x) = x - x^2/2$ corresponds to $\chi^2$ divergence. At the optimum, $r_Q$ separates immediate reward from continuation values via the Bellman constraint.

\subsection{GLADIUS \citep{kang2025}}

GLADIUS parameterizes both $Q_\theta(s,a)$ and $\hat{E}_\theta[V(s') \mid s,a]$ with neural networks, trained jointly on observed $(s,a,s')$ transitions:
\begin{equation}\label{eq:gladius}
\mathcal{L}(\theta) = -\frac{1}{N}\sum_{i}\log\pi_\theta(a_i \mid s_i)
  + \lambda\,\frac{1}{N}\sum_{i}\bigl(\hat{E}_\theta[V(s_i') \mid s_i,a_i] - V_\theta(s_i')\bigr)^{2},
\end{equation}
where $V_\theta(s') = \sigma\log\sum_{a'}\exp(Q_\theta(s',a')/\sigma)$. After training, reward is extracted via the same inverse Bellman operator as IQ-Learn: $\hat{r}(s,a) = Q_\theta(s,a) - \beta\,\hat{E}_\theta[V(s')\mid s,a]$. The MLP architecture shares parameters across states, providing implicit regularization over tabular $Q$.

%% ======================================================================
\section{Counterfactual Taxonomy}
%% ======================================================================

\begin{definition}[Type I: state extrapolation]
The transition kernel $P$ and reward $r$ are unchanged. Only the realized state values shift. The counterfactual CCPs require only evaluating the estimated Q at new state inputs, with no Bellman re-solve.
\end{definition}

\begin{definition}[Type II: environment change]
The reward $r$ is unchanged, but the transition kernel shifts to $\tilde{P}$ or the action set changes. The counterfactual Q must satisfy a new Bellman equation $\tilde{Q} = r + \beta \tilde{P}\tilde{V}$, which requires the structural reward $r$ separated from the old continuation values.
\end{definition}

\begin{proposition}[Type I equivalence]\label{prop:type1}
Under Type I, all four methods yield identical counterfactual CCPs.
\end{proposition}
\begin{proof}
Each recovers $Q^*(s,a) + c(s)$ for some action-independent $c$. The softmax $\pi(a \mid \tilde{s}) \propto \exp(Q^*(\tilde{s},a) + c(\tilde{s}))$ cancels $c$. \qed
\end{proof}

\begin{proposition}[Type II divergence]\label{prop:type2}
Under Type II, reduced-form produces invalid predictions. AIRL with anchors, IQ-Learn, and GLADIUS produce valid predictions when $r$ is correctly recovered.
\end{proposition}

The reduced-form estimate $Q^* + c(s)$ bundles $r$ with $\beta PV^*$. Under $\tilde{P}$, continuation values change but the reduced-form has no way to disentangle $r$. AIRL with anchors recovers $g^* = r^*$ exactly (Theorem~\ref{thm:recovery}). IQ-Learn and GLADIUS separate $r$ via the inverse Bellman operator. The LSW anchors provide what \citet{kalouptsidi2021} call the ``additional restrictions'' needed to identify counterfactuals under transition changes.

%% ======================================================================
\section{Application to Lee, Sudhir \& Wang (2026)}
%% ======================================================================

LSW study sequential content consumption with three actions (buy, wait, exit) using AIRL under anchors (A1)--(A2) and deterministic transitions. Their wait-time validation (Section 6.3) is a Type I counterfactual, so a reduced-form logit would pass this test identically. Their WFF removal and segment-specific delay policies (Sections 7.1--7.2) are Type II because they change the action-set structure and the MDP that each segment faces. Their content-based paywall (Section 7.3) is also Type II because high-content episodes become paid-only, altering the transition structure episode-by-episode.

%% ======================================================================
\section{Simulation Study}
%% ======================================================================

We construct a tabular MDP that mimics the LSW book-reading environment. The state encodes episode position, content quality, and wait-time level, giving $15 \times 3 \times 4 = 180$ regular states plus one absorbing state. The three actions are buy, wait, and exit, with deterministic transitions and $\beta = 0.95$. The true reward is $r^*(s,\text{buy}) = u(c) - 2.0$, $r^*(s,\text{wait}) = u(c) - 0.08w$, and $r^*(s,\text{exit}) = 0$, where $u \in \{0.5, 1.5, 3.0\}$. Analyses A through C and E through F are evaluated at the population level with no estimation noise. Analysis D uses finite-sample estimation.

\subsection{A. Reward recovery}

Methods with anchors pin the value function offset and recover $r^*$ exactly. Reduced-form logit converges to the advantage $A^* = r^* + \beta PV^* - V^*$, which is a noisy proxy for the reward because $V^*$ adds large state-specific offsets.

\begin{table}[H]
\centering
\begin{tabular}{lrr}
\toprule
Method & MSE vs.\ $r^*$ & Corr.\ with $r^*$ \\
\midrule
Oracle / AIRL+anchors / IQ / GLADIUS & 0.000 & 1.000 \\
AIRL-no-anchors & 10.005 & 0.440 \\
Reduced-form & 40.020 & 0.138 \\
\bottomrule
\end{tabular}
\caption{Population-level reward recovery. Type I CCP errors are below $10^{-6}$ for all methods.}
\label{tab:reward}
\end{table}

\subsection{B. Type II error vs.\ transition shift}

The counterfactual changes the transition kernel so that buying advances $k$ episodes instead of one. The Bellman equation is re-solved under each method's recovered reward.

\begin{table}[H]
\centering
\begin{tabular}{crrrr}
\toprule
$k$ & Oracle/AIRL+a/IQ/GLADIUS & AIRL-no-a & Reduced-form \\
\midrule
1  & 0.000 & 0.000 & 0.000 \\
2  & 0.000 & 0.025 & 0.082 \\
3  & 0.000 & 0.017 & 0.094 \\
5  & 0.000 & 0.015 & 0.123 \\
7  & 0.000 & 0.016 & 0.129 \\
10 & 0.000 & 0.014 & 0.139 \\
\bottomrule
\end{tabular}
\caption{Type II CCP error. Structural methods are exact. Reduced-form error grows with shift magnitude.}
\label{tab:type2}
\end{table}

\subsection{C. Shaping magnitude sweep}

We construct shaped rewards $\tilde{r} = r^* + \alpha V^*(s) - \alpha\beta V^*(s')$ for $\alpha \in [0,1]$ and evaluate the Type II counterfactual with $k=3$.

\begin{table}[H]
\centering
\begin{tabular}{rrr}
\toprule
$\alpha$ & Type II CF error & Corr.\ with $r^*$ \\
\midrule
0.00 & 0.000 & 1.000 \\
0.10 & 0.004 & 0.908 \\
0.30 & 0.011 & 0.610 \\
0.50 & 0.017 & 0.440 \\
1.00 & 0.027 & 0.270 \\
\bottomrule
\end{tabular}
\caption{CF error is monotone in shaping magnitude. Without anchors, $\alpha$ is indeterminate.}
\label{tab:shaping}
\end{table}

\subsection{D. Finite-sample convergence}

We compare three finite-sample estimators across sample sizes ranging from 200 to 10{,}000 trajectories (3 seeds each). The reduced-form logit converges to $A^*$, which is the wrong target for Type II counterfactuals. IQ-Learn ($\chi^2$) uses the actual tabular implementation from \citet{garg2021} with L-BFGS-B optimization. GLADIUS uses neural $Q$ and EV networks with 150 training epochs.

\begin{table}[H]
\centering
\begin{tabular}{rrrr}
\toprule
$N$ & Reduced-form & IQ-Learn ($\chi^2$) & GLADIUS \\
\midrule
200      & 0.265 & 0.237 & 0.132 \\
500      & 0.264 & 0.287 & 0.135 \\
2{,}000  & 0.263 & 0.285 & 0.139 \\
5{,}000  & 0.262 & 0.335 & 0.135 \\
10{,}000 & 0.262 & 0.333 & 0.130 \\
\bottomrule
\end{tabular}
\caption{Type II CF error ($k=3$, 3 seeds). GLADIUS cuts error by approximately 50\% over tabular methods via neural regularization. IQ-Learn ($\chi^2$) does not improve over reduced-form logit because the tabular chi-squared objective encounters poor local optima as $N$ grows and the loss surface sharpens.}
\label{tab:samplesize}
\end{table}

\subsection{E. Anchor misspecification}

We set the true exit payoff to $\varepsilon$ while the anchors assume it is zero.

\begin{table}[H]
\centering
\begin{tabular}{rrrr}
\toprule
$\varepsilon$ & Corr.\ with $r^*$ & Type II CF error & In-sample $\pi$ shift \\
\midrule
0.0 & 1.000 & 0.000 & 0.000 \\
0.5 & 0.976 & 0.002 & 0.000 \\
1.0 & 0.909 & 0.004 & 0.000 \\
2.0 & 0.729 & 0.008 & 0.000 \\
\bottomrule
\end{tabular}
\caption{Degradation is smooth, not knife-edge. In-sample CCPs are unaffected.}
\label{tab:misspec}
\end{table}

\subsection{F. Content-based paywall}

High-content episodes become paid-only and the wait action is blocked, forcing buy-or-exit. This is a Type II shift that mirrors Section 7.3 of LSW.

\begin{table}[H]
\centering
\begin{tabular}{lr}
\toprule
Method & CCP error \\
\midrule
Oracle / AIRL+anchors / IQ-Learn & 0.000 \\
AIRL-no-anchors & 0.009 \\
Reduced-form & 0.185 \\
\bottomrule
\end{tabular}
\caption{Reduced-form fails at 18.5pp because it cannot separate content utility from continuation value.}
\label{tab:paywall}
\end{table}

%% ======================================================================
\section{Summary}
%% ======================================================================

\begin{table}[H]
\centering
\begin{tabular}{lcccc}
\toprule
 & Reduced-form & AIRL (generic) & AIRL (anchors) & IQ-Learn / GLADIUS \\
\midrule
Recovers         & $A^*(s,a)$ & $r$ + shaping & $r^*(s,a)$ & $r^*(s,a)$ \\
Type I CF        & \checkmark & \checkmark    & \checkmark & \checkmark \\
Type II CF       & $\times$   & Only if $\delta \equiv 0$ & \checkmark & \checkmark \\
\bottomrule
\end{tabular}
\label{tab:summary}
\end{table}

%% ======================================================================

\begin{thebibliography}{99}

\bibitem[Fu et~al.(2018)]{fu2018}
Fu, J., Luo, K., and Levine, S. (2018). ``Learning Robust Rewards with Adversarial Inverse Reinforcement Learning.'' ICLR.

\bibitem[Garg et~al.(2021)]{garg2021}
Garg, D., Chakraborty, S., Cundy, C., Song, J., and Ermon, S. (2021). ``IQ-Learn: Inverse soft-Q Learning for Imitation.'' NeurIPS.

\bibitem[Kalouptsidi et~al.(2021)]{kalouptsidi2021}
Kalouptsidi, M., Scott, P.~T., and Souza-Rodrigues, E. (2021). ``Identification of Counterfactuals in Dynamic Discrete Choice Models.'' Quantitative Economics, 12(2), 351--403.

\bibitem[Kang et~al.(2025)]{kang2025}
Kang, E.~H., Yoganarasimhan, H., and Jain, L. (2025). ``An Empirical Risk Minimization Approach for Offline Inverse RL and Dynamic Discrete Choice Models.'' arXiv:2502.14131.

\bibitem[Lee et~al.(2026)]{lee2026}
Lee, P., Sudhir, K., and Wang, T. (2026). ``Consumer Engagement with Sequential Content: A Content-Aware Dynamic Choice Model.'' Working paper, Yale University.

\bibitem[Magnac and Thesmar(2002)]{magnac2002}
Magnac, T. and Thesmar, D. (2002). ``Identifying Dynamic Discrete Decision Processes.'' Econometrica, 70(2), 801--816.

\bibitem[Ng et~al.(1999)]{ng1999}
Ng, A.~Y., Haley, D., and Russell, S.~J. (1999). ``Policy Invariance Under Reward Transformations.'' ICML.

\bibitem[Rust(1987)]{rust1987}
Rust, J. (1987). ``Optimal Replacement of GMC Bus Engines.'' Econometrica, 55(5), 999--1033.

\end{thebibliography}

\end{document}
